\chapter{Survey Instrument and Client Meeting Evidence}
\label{app:survey}

This appendix documents the empirical basis of the requirements analysis
(Chapter~\ref{cap:analysis}) and the validation (Chapter~\ref{cap:validation}).
Section~\ref{app:survey-meeting} summarises the client engagement that
produced the pain points; Section~\ref{app:survey-instrument} reproduces the
full survey instrument used to validate them.

\section{Client Engagement and Scrum Cadence}
\label{app:survey-meeting}

The project was conducted as a five-month internship (16~February~--~9~July
2026) with adidas Business Services in Porto, a shared-services centre
employing more than 900 people, of whom roughly 22\% are international,
operating daily across some 15 business languages. The original client brief
asked for a \emph{framework and sustainable strategy} to attract, develop and
retain language capabilities; the platform documented in this report is the
concrete proof-of-concept that operationalises that brief.

\begin{table}[ht]
	\centering
	\caption{Selected milestones of the client engagement.}
	\label{tab:meeting-log}
	\begin{tabular}{@{}l p{8.5cm}@{}}
		\toprule
		\textbf{Date} & \textbf{Milestone} \\
		\midrule
		Feb 2026      & Problem analysis and requirements gathering; original four-question survey bank drafted. \\
		Late Feb 2026 & Kick-off meeting (Hellenic Mediterranean University, Crete) with the international, multidisciplinary team. \\
		Feb--Mar 2026 & System architecture and database design. \\
		Mar--Apr 2026 & Core platform development: CV parsing, candidate management, job matching. \\
		Apr 2026      & Language-assessment module integration. \\
		29 Apr 2026   & Client pain-point review; nine pain points consolidated and mapped to survey questions. \\
		Apr--May 2026 & Testing, deployment and stakeholder validation on the production environment. \\
		May 2026      & Survey distribution, data collection and analysis. \\
		May--Jun 2026 & Final report and presentation. \\
		\bottomrule
	\end{tabular}
\end{table}

These meetings established the Scrum-with-the-client cadence described in
Section~\ref{sec:method-scrum}: the author acted as Product Owner, translating
client conversations into a prioritised backlog and demonstrating increments
on the deployed environment. Nine pain points emerged from the engagement and
are reproduced verbatim below; each maps to one or more platform modules in
Table~\ref{tab:survey-mapping}.

\begin{enumerate}
	\item \textbf{Organise candidate information} --- the team sometimes relies on Excel.
	\item \textbf{Know which candidates applied before or are new} --- under GDPR, candidate data may be kept for only six months.
	\item \textbf{Screening for language skills} (extendable to any skill) --- a first screening layer, not a replacement for the existing process.
	\item \textbf{Direct contact with candidates} --- the basis of the notifications and campaigns features.
	\item \textbf{Save time analysing CVs and extracting information} --- even reading a candidate's location requires opening every CV individually.
	\item \textbf{Analytics} --- clear, organised views of the recruitment process.
	\item \textbf{Understand where to focus recruitment efforts} --- e.g. German speakers, internship candidates, technical profiles.
	\item \textbf{Empower HR productivity} --- a tool that makes the work feel more structured.
	\item \textbf{Keep non-selected candidates engaged} --- re-engaging ``borderline'' candidates within the six-month window via assessments, campaigns and improvement tracks.
\end{enumerate}

\section{Survey Instrument}
\label{app:survey-instrument}

The survey was designed so that each question describes a problem neutrally,
without hinting at any solution, so that responses constitute independent
evidence for the corresponding module. One question was authored per pain
point, with backup questions where a single phrasing was insufficient.

\subsection*{Pain-point questionnaire}

\begin{enumerate}
	\item \textit{Organising candidate information.}\\
	``What tool does your team primarily use to store and manage candidate
	information during a recruitment process?''
	(Excel / Google Sheets / An ATS / Email folders / Shared documents /
	No structured system)

	\item \textit{Knowing if a candidate has applied before (GDPR / six-month rule).}\\
	``How does your team currently identify whether a candidate has previously
	applied for a position at adidas Porto?''
	(Check manually / Rely on the candidate declaring it / No way to verify /
	Automated system) \\
	\emph{Backup:} ``Has your team ever lost useful information about a past
	candidate because data-retention policies required it to be deleted?''

	\item \textit{Language / skills screening.}\\
	``At what stage of your recruitment process do you currently perform a
	first formal check of a candidate's language proficiency or technical
	skills?''
	(Before the first interview / During the first interview / Only in a
	dedicated assessment round / No formal check --- we rely on the CV)

	\item \textit{Direct contact and communication.}\\
	``How does your team currently keep candidates informed about their
	application status or upcoming opportunities?''
	(Manual individual emails / Automated emails via an ATS / Phone calls /
	No structured process) \\
	\emph{Backup:} ``Has your team ever had to re-engage a candidate for a new
	position after they were not selected for a previous one?''

	\item \textit{Saving time on CV analysis.}\\
	``When screening a batch of CVs, how long does it typically take your team
	to extract key information (e.g.\ location, years of experience, languages)
	from each individual CV?''
	(Under 1~min / 1--3~min / 3--5~min / More than 5~min) \\
	\emph{Backup:} ``How often do you open a CV only to find immediately that
	the candidate does not meet a basic requirement?''

	\item \textit{Analytics and pipeline visibility.}\\
	``How easy is it for you to get a quick overview of your current
	recruitment pipeline --- how many candidates are at each stage, or how
	many applications were received per opening?''
	(Very easy --- I have a dashboard / Possible but slow / Very difficult /
	No visibility)

	\item \textit{Where to focus recruitment efforts.}\\
	``When planning a new recruitment campaign, how does your team identify
	which skills, languages or profiles are most in demand?''
	(Analyse past data / Rely on managers / Done intuitively / No structured
	process)

	\item \textit{HR tooling and productivity.}\\
	``On a scale of 1 to 5, how would you rate the degree to which your current
	recruitment tools help you work more efficiently and reduce repetitive
	manual tasks?''
	(1 = very little value $\to$ 5 = significant productivity boost)

	\item \textit{Keeping borderline candidates engaged.}\\
	``What happens to candidates who were not selected for a position but
	showed potential --- for example strong language skills or relevant
	experience?''
	(Keep their profile for future openings / Inform them proactively /
	Rejected and removed / No structured process)
\end{enumerate}

\subsection*{Original four-question bank (28 April session)}

A shorter instrument, drafted before the full pain-point mapping, framed four
quantitative questions later folded into the questionnaire above:

\begin{enumerate}
	\item \textit{CV screening workload} --- ``On average, how many hours per
	week does your team spend manually reviewing CVs and shortlisting
	candidates for a single open position?'' ($<$1h / 1--3h / 3--6h / $>$6h)
	\item \textit{Language proficiency verification} --- ``How confident are
	you in your current process for objectively verifying a candidate's foreign
	language proficiency before the first interview?'' (1--5)
	\item \textit{Evaluation consistency across the team} --- ``When multiple
	HR members evaluate the same candidate, how often do inconsistencies arise
	in how they are scored or ranked?'' (Never $\to$ Always)
	\item \textit{Tracking candidate communication history} --- ``How easy is
	it to retrieve a complete history of all interactions with a specific
	candidate across different HR team members?'' (1--5)
\end{enumerate}

\subsection*{Response data}

The questionnaire was completed in May~2026 by the lead HR stakeholder at
adidas Porto, the same person who framed the original client brief. Because
the respondent speaks for the HR function that owns the problem, the response
is treated as a structured validation interview rather than a statistical
sample; its analysis and the three headline findings are discussed in
Section~\ref{sec:val-survey}. Table~\ref{tab:survey-responses} reproduces the
selected answers verbatim. Notably, the respondent's team already operates an
enterprise Applicant Tracking System (SAP SuccessFactors), which makes the
residual friction recorded below particularly significant.

\begin{longtable}{@{}r p{6.0cm} p{6.0cm}@{}}
	\caption{Client responses to the validation survey (May 2026).}
	\label{tab:survey-responses} \\
	\toprule
	\textbf{\#} & \textbf{Question (abbreviated)} & \textbf{Selected response} \\
	\midrule
	\endfirsthead
	\multicolumn{3}{c}{\tablename\ \thetable\ -- \textit{continued}} \\
	\toprule
	\textbf{\#} & \textbf{Question (abbreviated)} & \textbf{Selected response} \\
	\midrule
	\endhead
	\midrule \multicolumn{3}{r}{\textit{continued on next page}} \\
	\endfoot
	\bottomrule
	\endlastfoot

	1  & Primary tool to store/manage candidate information & \textit{Other:} SAP SuccessFactors (enterprise ATS) \\
	2  & How repeat applicants are identified & We use an automated system \\
	3  & Ever lost useful candidate data to retention rules & Yes, occasionally \\
	4  & Stage of first formal language / skill check & Before the first interview; and only in a dedicated assessment round \\
	5  & How candidates are kept informed & Individual manual emails; automated ATS emails; phone calls \\
	6  & Ease of re-engaging a previously rejected candidate & Yes, and it was easy \\
	7  & Time to extract key fields from one CV & Under 1 minute \\
	8  & Frequency of opening a clearly unsuitable CV & Rarely \\
	9  & Ease of getting a pipeline overview & \textit{Other:} Very easy, using the ATS \\
	10 & How in-demand skills/profiles are identified & We rely on input from managers \\
	11 & Share of CVs that clearly fail minimum requirements & 30--60\% \\
	12 & Rating (1--5) of how much current tools boost productivity & \textbf{1 --- they add very little value} \\
	13 & What happens to promising but non-selected candidates & Keep their profile for future openings; inform them proactively --- \textit{Other:} ``the process here has still space for improvement'' \\
	14 & Ease of knowing what actions colleagues already took & \textit{Other:} ``Sometimes difficult --- you need to enter each process or ask the colleague'' \\

\end{longtable}

The three responses with the greatest validating weight are highlighted in
Section~\ref{sec:val-survey}: the floor productivity rating (Q12) from a
team that already runs an enterprise ATS; the estimate that 30--60\% of
received CVs fail the minimum requirements (Q11); and the reliance on
managerial intuition rather than data for strategic skill planning (Q10).
Together with the cross-team coordination difficulty (Q14) and the
acknowledged room for improvement in candidate re-engagement (Q13), they
provide direct, source-authoritative evidence that the platform's
prioritised feature set addresses real operational needs.
